\section{Philosophy, Poetry, Piety}

\paragraph{Philosopher}


The \textbf{Philosopher}, in the Traditional sense, is a lover of Wisdom, which is the intuitive knowledge of the Ideas. The goal of the Philosopher is, then, to be a Wise Man, or Sage. His method is the examined life, or self-knowledge, which differs from the natural knowledge of the scientist.

This knowledge leads him to understand the best way to live a human life, specifically, knowledge of the virtues. On the metaphysical principle that knowing is equivalent to being, the Philosopher can know the virtues only by becoming virtuous, not through definitions or sophistical arguments. Of course, by ``virtue'', we mean a strength or the power to act appropriately in a given situation.

Justice, wisdom, courage, temperance are the cardinal virtues, each of the latter of which apply predominantly to a different class, balanced by justice. Of course, this list of virtues is not exhaustive. By seeking to actualize all the virtues, that is, all of his possibilities, the Philosopher is more like the True Man\footnote{\url{http://www.gornahoor.net/?p=1039}}. Since the True Man is characteristic of the Kshatriya caste, only the Philosopher or Sage is entitled to rule the Polis. Before the objection is raised that the philosopher is impractical, let us remember that Socrates had proved his courage in battle.

However, the common man is ruled by passion rather than reason, so he has no interest in philosophy, nor can he understand it. Hence, the common man rejects the rule of the Sage. This is the paradoxical position that the Sage finds himself in: he is the only one capable of ruling, but he is not able to rule.

\paragraph{Poet}


The \textbf{Poet}, in the larger sense, include the poet properly so-called, the playwright, the singer, the lyrical musician, and in modern times, the actor and the filmmaker. In short, any art that appeals to the ``right brain'' in popular parlance, or more accurately the bicameral mind as described by Julian Jaynes\footnote{\url{https://www.gornahoor.net/?p=1365}}.

The common man learns the virtues from the Poet, not the Philosopher. Thus poems, songs, myths, stories, and film are the vehicles through which the virtues and the good life are taught to the masses. Unfortunately, in our times, the ideas of the virtues and good life are defective, since they are all too often devoid of any influence of Wisdom, or else they are one-dimensional. Nevertheless, there have been many examples of Poets and poetry that show understanding of the virtues.

Paradoxically, the poet need not be virtuous in order to portray the virtues, he is instead the imitator of the virtues. Thus, a craven actor may nevertheless portray a great hero in a play or film. A bookish novelist may write deep and perceptive speeches for his men-of-the-world characters. This is the illusion of art, since there are no bad consequences. A play can bomb, but the playwright lives to create another despite his hurt feelings. But, in battle, a bad decision can be fatal, not just to the leader, but potentially to the nation he pledged to defend.

Nevertheless, the Poet must have some knowledge of the virtues, though not like the Philosopher's knowledge. It is an infused knowing, like an unconscious act of grace, a form of madness, a gift from the Muses\footnote{\url{https://www.gornahoor.net/?p=1613}}.

\paragraph{Pious}


\textbf{Piety}, or the worship of the ancestral gods, is also a virtue.

But we know about the gods from the Poets, not the Philosophers, as their acts are passed on through myths. This puts the Philosopher in an awkward position. On the one hand, he desires the good life for his family and nation. Thus he wants to preserve the Traditions of his ancestors and the way of life of his people, since they have served as the bases of public order and the common good.

On the other hand, his understanding of the gods is not gained through myths, which mediate knowledge through sensual imagery. Rather, it derives from an unmediated intuition of the Ideas. This intuition does not involve any of the external senses, not even in the imagination. Yet, there is a profound relationship between this intuition and its mythological representation, as described in The Word is Silence Expressed\footnote{\url{https://www.meditationsonthetarot.com/the-word-is-silence-expressed}}.


\flrightit{Posted on 2011-02-15 by Cologero}

\begin{center}* * *\end{center}

\begin{footnotesize}
\begin{sffamily}
\texttt{Cologero on 2011-02-19 at 10:33 said:}

The discussion on the cardinal virtues again illustrates the continuity from the ancient Greek tradition to the ancient Roman tradition to the period of Christendom. The virtues must be understood in the same way from age to age.

If those virtues have been replaced by other so-called virtues, such as ``tolerance'' or universalism and so on, that is the result of deviations of the modern age, not of the alleged return to some ``pure'' or ``primitive'' Christianity (which never existed as such in the first place.)


\hfill

\end{sffamily}
\end{footnotesize}

